\section{Peak-to-Total-Verhältnisse}

Es wird das Peak\-to\-Total\-Verhältnis (PTT) für die Messungen beider Detektoren von Cäsium und Cobalt bestimmt. Die Fläche der Peaks lässt sich Tabelle \ref{tab:nai_co_peaks}, \ref{tab:nai_cs_peaks}, \ref{tab:ge_co_peaks} und \ref{tab:ge_cs_peaks} entnehmen, wobei die Zuordnung der Peaks zu den Energien in den Tabellen \ref{tab:zuordnungen_ge} und \ref{tab:zuordnungen_nai} oder über die Energiekalibration möglich ist.

Die Gesamtanzahl der registrierten Photonen wird über das Summieren der Spektren abzüglich der Untergrundmessungen ermittelt und ist in Tabelle \ref{tab:sum_counts} zu finden. Bei Cäsium werden die beiden Peaks addiert. Es ist also $PTT = A / \Sigma_{i=0}^{16383} n_i$.Damit erhält man die PTT in Tabell \ref{tab:ptts}.


\begin{table}[h]
    \centering
    \begin{tabular}{ccc}
	& Cäsium & Cobalt \\
	NaI & $\num{0.3503(11)}$ & $\num{0.248(11)}$\\
	HPGe & $\num{0.22854(20)}$ & $\num{0.14420(28)}$\\
    \end{tabular}
    \caption{Berechnete PTT}
    \label{tab:ptts}
\end{table}


Wie zu erwarten sind die PTT im Szintillator größer als die im HPGe Detektor. Auch nehmen die PTT bei zunehmenden Photonenenergien ab. Dies passt durchweg gut zu den Beobachtungen aus Abschnitt \ref{sec:efficiency}.
